\section{Introduction}
\label{sec:intro}

\subsection{Motivation}

Multi-agent language-model systems are now a common deployment pattern. A single model
is instantiated several times, the instances answer or debate, and their outputs are
combined by voting or by an orchestrator. The empirical case for this pattern rests on
accuracy. \citet{du2023debate} showed that multi-agent debate improves factuality and
arithmetic reasoning relative to a single pass, and a large body of subsequent work
adopted debate and orchestration as general-purpose accuracy levers.
\citet{kim2025scaling} argue for a quantitative science of agent-system scaling and
propose coordination efficiency, communication overhead, error amplification and
redundancy as the relevant system-level quantities.

The optimistic reading is contested. \citet{smit2024debate} find that debate protocols
are highly sensitive to cost and hyperparameter choices and are not reliably superior to
self-consistency at matched compute. \citet{wynn2025debate} document failure modes in
which communication moves agents toward agreement on incorrect answers. Both critiques
are about accuracy. Neither addresses the question that matters once such a system is
deployed behind an interface that reports how confident it is.

That question is calibration. A deployed multi-agent system almost always exposes some
confidence signal: the fraction of agents that agreed, the orchestrator's own
probability, or a verbalized percentage. Calibration of single models is well studied.
\citet{guo2017calibration} established that modern neural networks are systematically
overconfident and that expected calibration error is the standard measurement. What is
not established is what coordination does to calibration. Aggregating several agents
changes both the answer and the confidence signal attached to it, and there is no reason
in principle for the aggregate signal to inherit the calibration of its parts.

The two existing positions disagree. \citet{yang2024collaborative} report that
deliberation among agents improves confidence assessment, which implies coordination is
a calibration benefit. \citet{wang2024calibrating} show that self-consistency, a
degenerate form of aggregation over samples rather than agents, improves calibration on
mathematical reasoning. Against this, if debate induces correlated errors as
\citet{wynn2025debate} describe, then agreement becomes a less honest signal precisely
when the team is wrong, and coordination is a calibration cost.

A third consideration cuts across both. Test-time reasoning budgets are now a first-class
scaling axis. \citet{zhu2025testtime} show that additional inference compute improves
agent performance, and \citet{li2026agentbench} benchmark test-time scaling for general
agents and note that sequential scaling eventually meets context ceilings. Whether
reasoning tokens substitute for parameters, and whether they help or harm calibration,
is not settled. \citet{wong2026teamofthoughts} propose orchestrated tool calling as an
efficient test-time scaling mechanism, which presupposes that orchestration is worth its
overhead.

\subsection{What this sweep measures}

This report documents a factorial sweep designed to separate these effects. The design
crosses four continuous or ordinal levers against four coordination structures, four
task families and three replication seeds, holding the model family, the serving stack,
the question set and the grading procedure fixed. The primary outcome is not accuracy
alone but the pair (accuracy, calibration), and the primary estimand is

\begin{equation}
\dece \;=\; \mathrm{ECE}\bigl(c_{\mathrm{system}}\bigr) \;-\; \overline{\mathrm{ECE}\bigl(c_{\mathrm{agent}}\bigr)},
\label{eq:estimand}
\end{equation}

the difference between the calibration error of the confidence signal a system exposes
and the mean calibration error of the individual agents inside it. A positive value means
coordination degraded calibration relative to the parts it was built from; a negative
value means coordination corrected it.

Two system signals are measured separately, and the report keeps them separate
throughout, because they answer different design questions. The vote-fraction signal,
$\dvote$, asks whether the team's level of agreement is an honest probability. The
final-producer signal, $\dfp$, asks whether the agent whose output actually determines
the answer remains well calibrated after coordination. A product that displays ``three
of three agents agree'' is exposing the first; a product that displays a calibrated model
probability is exposing the second.

The models are open weights, which is what makes the estimand measurable at all. Option
logprobs, realized thinking-token counts, per-round per-agent outputs and the exact
prompt token identities are all recorded, so the per-agent baseline in
Equation~\ref{eq:estimand} is observed rather than inferred.

\subsection{Scope and structure}

The sweep reported here contrasts one-agent and three-agent systems. Agent count is not
varied continuously: every multi-agent cell uses three agents, so the one-against-three
comparison is a comparison of system families rather than a point on an agent-count
scaling curve. Sibling runs with two, four, five, six and seven agents exist as frozen
cell manifests but have not been ingested, and are discussed only as future work in
Section~\ref{sec:limitations}.

Section~\ref{sec:design} states the design: the axes, their levels, the factorial size,
the trim that reduced it, and the collapsing rules that make certain factor combinations
degenerate. Section~\ref{sec:system} documents the implementation: coordination
protocols, context-sharing rendering, prompt templates, the reasoning-budget mechanism,
the serving stack and the benchmark contracts. Section~\ref{sec:measurement} documents
measurement and inference: confidence elicitation, the calibration estimators, the
efficiency metrics and the two bootstrap schemes.

Sections~\ref{sec:health} through~\ref{sec:items} report results. They are deliberately
exhaustive: every quantity the analysis pipeline computed appears in the text, in a
table, or in a figure, and in most cases in all three. Section~\ref{sec:interpretation}
assesses which findings are novel relative to the literature cited above and what would
be required to establish them. Section~\ref{sec:limitations} states the threats to
validity. Appendices~\ref{app:axes} through~\ref{app:provenance} carry the complete
tables, the full figure atlas and the provenance record.
